% !TEX root = main.tex

\documentclass{article}
\usepackage[T1]{fontenc}
\usepackage[utf8]{inputenc}
\usepackage{amsmath}
\usepackage{amssymb}
\usepackage{hyperref}
\usepackage{parskip} %skip the indent of a new paragraph.
\usepackage{float}
\usepackage{graphicx}
\usepackage{listings}
\usepackage[per-mode=symbol]{siunitx}
\usepackage{epstopdf}
\usepackage[super]{nth}
\lstset{language=Matlab, frame=single, breaklines=true,numbers=left, keywordstyle=\color{blue},rulecolor=\color{black},commentstyle=\color{gray}}

\usepackage{cleveref}
\usepackage{todonotes}

%% Brukes for tabeller (av likninger)
\usepackage{tabularx}
\def\tabularxcolumn#1{m{#1}}

\newcommand{\mbf}[1]{\mathbf[#1]}
\newcommand{\partialderiv}[2]{\frac{\partial{#1}}{\partial{#2}}} % prints partial derivative as a fraction

\begin{document}
    % !TEX root = main.tex

\begin{titlepage}
    %\maketitle
    %\rule{\linewidth}{0.5mm}
    \begin{center}
    	\large
    	Optimization \& Control TTK4135
    \end{center}
    \vspace{\fill}
    \rule{\linewidth}{0.5mm}
    \begin{center}
    	\huge
    	Helicopter Lab
    \end{center}
	\rule{\linewidth}{0.5mm}
	\vspace{\fill}

	%\begin{center}
    %	\huge
    %	Bern Johan Damslora -- 759477 \\ Didrik Rokhaug -- 759528
    %\end{center}
    \large
    \centering
    \begin{table}[H]
    	\centering
    	\large
    	\begin{tabular}{c}
            nr. 999999\\
    		nr. 759528\\
            nr. 999999
    	\end{tabular}
    \end{table}
    \vspace{\fill}
    \begin{center}
    	\large
    	\today
    \end{center}
	\vspace{\fill}
    \begin{figure}[H]
    \centering
    \includegraphics[width=0.5\textwidth]{logontnu_eng}
    \end{figure}
    \thispagestyle{empty}
\end{titlepage}

    % !TEX root = main.tex

\section{Optimal control of Pitch/Travel without Feedback}
\subsection{1}
The continuous time system can be written as $\mathbf{\dot{x}} = \mathbf{A}_c\mathbf{x}+ \mathbf{B}_c\mathbf{u}$, with
\begin{subequations}
    \begin{align}
        \mathbf{x} &= \begin{bmatrix}
            \lambda\\
            r\\
            p\\
            \dot{p}
        \end{bmatrix}\\
        \mathbf{A}_c &= \begin{bmatrix}
            0 & 1 & 0 & 0\\
            0 & 0 & -K_2 & 0\\
            0 & 0 & 0 & 1\\
            0 & 0& -K_1K_{pp} & -K_1K_{pd}
        \end{bmatrix} \\
        \mathbf{B}_c &= \begin{bmatrix}
            0\\
            0\\
            0\\
            K_1K_{pp}
        \end{bmatrix}
    \end{align}
\end{subequations}

This is a model of the helicopter with pitch controller. The model inputs are set points for the pitch controller.

\subsection{2}
The system was discretized using forward Euler:
\begin{equation}
    \mathbf{\dot{x}} = \mathbf{A}_c\mathbf{x}_k + \mathbf{B}_c\mathbf{u}_k \approx \frac{\mathbf{x}_{k+1}-\mathbf{x}}{h}
\end{equation}

Rearranging, we get
\begin{equation}
    \mathbf{x}_{k+1} = (\mathbf{A}_ch + \mathbf{I})\mathbf{x}_k + h\mathbf{B}_c\mathbf{u}_k
\end{equation}

and we see that
\begin{subequations}
    \begin{align}
        \mathbf{A}_d &= \mathbf{A}_ch+\mathbf{I} \\
        \mathbf{B}_d &= h\mathbf{B}_c
    \end{align}

\end{subequations}

\begin{subequations}
    \begin{align}
        \mathbf{A_d} &= \begin{bmatrix}
        0 & 1 & 0 & 0\\
        0 & 0 & -K_2 & 0\\
        0 & 0 & 0 & 1\\
        0 & 0 & -K_1K_{pp} & -K_1K_{pd}
        \end{bmatrix}\\
        \mathbf{B_d} &= \begin{bmatrix}
        0\\
        0\\
        0\\
        K_1K_{pp}
    \end{bmatrix}
    \end{align}
\end{subequations}

\subsection{3}
The optimal trajectory was calculated using quadprog, which takes a weight matrix Q, vector c (from the cost function $x^TQx + c^Tx$), inequality constraints matrices A, B, equality constraints Aeq, Beq and lower and upper bounds vlb, vub on the inputs, and returns a vector z with the optimal trajectory for all the states, along with a vector \lambda containing the Lagrange multipliers at the solution.

\begin{lstlisting}
    [z,lambda] = quadprog(Q,c,[],[], Aeq, beq, vlb, vub);
\end{lstlisting}

\begin{lstlisting}
    Q1(1,1) = 1;
    ...
    P1 = 10;  
    Q = 2*genq2(Q1,P1,N,M,mu);
    c = zeros(N*mx+M*mu,1); 
    ...
    Aeq = gena2(A1,B1,N,mx,mu);         
    beq = zeros(size(Aeq,1),1);
    ...
    [vlb,vub]       = genbegr2(N,M,xl,xu,ul,uu); 
\end{lstlisting}
The weight matrix Q was generated using the function genq2, and c is a vector of zeros. Aeq was generated using the function gena2, while beq is a vector of zeros. The constraints on the inputs were implemented using the given function genbegr2. 

Finally, mx denotes the number of states, mu the number of inputs, N the time horizon for states and M the time horizon for inputs. For a more detailed description of the above functions, see Appendix X. 

The cost function \todo{insert cost function (15)} consists of two terms, $\lambda_i - \lambda_f)^2$ and $qp_{ci}^2$. The first one aims to penalize deviation from the reference point $\lambda_f$, while the second one penalizes inputs $p_{ci}^2$. The relative cost of the two terms is set using the weight q. This is reflected 

One possible downside to the cost term $(\lambda_i - \lambda_f)^2$ could be that it does not take into account that $\lambda_i = \lambda_ i + 2\pi$, resulting in very poor pathing if exposed to a large perturbation. On the other hand, it might be desirable to actually do one or multiple revolutions.
\todo{finner vi flere ulemper?}

\subsection{4}
\todo{insert code + plots}
The helicopter does not end in the desired point $x_f$ for any of the states. This is largely \todo{er det flere grunner?} because of model inaccuracies; discretization errors from using forward Euler and linearization errors, in addition to inherent model errors/simplifications. Also, the system is affected by noise (e.g. wind currents) as well as measurement errors.

    % !TEX root = main.tex

\section{Optimal Control of Pitch/Travel with Feedback (LQ)}
\subsection{3}
MPC can be realized by solving a finite horizon open loop optimal control problem at each time step, \todo{how to do dis in simulink}

One advantage with MPC compared to LQR is that MPC is more robust. This is because the MPC calculates a new optimal input for each timestep, rather than doing everything a priori. In addition, MPC allows the use of constraints.

\missingfigure{\label{fig:prob3_lambda}}

On the other hand, MPC is far more resource-demanding than LQR\todo{LQ eller LQR?}. Also, the solution using MPC is usually not optimal in the context of the cost function (as mentioned above, it optimizes with respect to stability or other metrics). Another downside to MPC is that it requires a model of the system, while LQR does not.

    % !TEX root = main.tex

\section{Exercise 4}
\subsection{Continuous time state space model}
We introduce two new states: elevation $e$ and its rate $\dot{e}$. The continuous time system can now be written as $\mathbf{\dot{x}} = \mathbf{A}_c\mathbf{x}+ \mathbf{B}_c \mathbf{u}$, with
\begin{subequations}
    \begin{align}
        \mathbf{x} &= \begin{bmatrix}
            \lambda\\
            r\\
            p\\
            \dot{p}\\
            e\\
            \dot{e}\\
        \end{bmatrix}\\
        \mathbf{u} &= \begin{bmatrix}
            p_c\\
            e_c\\
        \end{bmatrix}\\
    \mathbf{A_c} &= \begin{bmatrix}
        0 & 1 & 0          & 0          & 0          & 0         \\
        0 & 0 & -K_2       & 0          & 0          & 0         \\
        0 & 0 & 0          & 1          & 0          & 0         \\
        0 & 0 & -K_1K_{pp} & -K_1K_{pd} & 0          & 0         \\
        0 & 0 & 0          & 0          & 0          & 1         \\
        0 & 0 & 0          & 0          & -K_3K_{ep} & -K_3K_{ed}
    \end{bmatrix}\\
    \mathbf{B_c} &= \begin{bmatrix}
        0         & 0        \\
        0         & 0        \\
        0         & 0        \\
        K_1K_{pp} & 0        \\
        0         & 0        \\
        0         & K_3K_{ep}
    \end{bmatrix}
    \end{align}
\end{subequations}
\subsection{Discrete time state space model}
The system was discretized using the forward Euler method. Using \cref{eq:eulerfwd}, we get the model $\mathbf{x_{k+1}} = \mathbf{A}_d\mathbf{x_k}+\mathbf{B}_d \mathbf{u_k}$, where
\begin{subequations}
    \begin{align}
    \mathbf{A_d} &= \begin{bmatrix}
        1 & h & 0           & 0            & 0           & 0         \\
        0 & 1 & -hK_2       & 0            & 0           & 0         \\
        0 & 0 & 1           & h            & 0           & 0         \\
        0 & 0 & -hK_1K_{pp} & 1-hK_1K_{pd} & 0           & 0         \\
        0 & 0 & 0           & 0            & 1           & h         \\
        0 & 0 & 0           & 0            & -hK_3K_{ep} & 1-hK_3K_{ed}
    \end{bmatrix}\\
    \mathbf{B_d} &= \begin{bmatrix}
        0          & 0        \\
        0          & 0        \\
        0          & 0        \\
        hK_1K_{pp} & 0        \\
        0          & 0        \\
        0          & hK_3K_{ep}
    \end{bmatrix}
    \end{align}
\end{subequations}

\subsection{Objective function and inequality constraints on elevation}
In this exercise, we minimized the cost function
\begin{equation}
        \phi =\sum (\lambda_i - \lambda_f)^2 + r_{p}p_{ci}^2 + r_{e}e_{ci}^2
\end{equation}
which can be written as
\begin{subequations}
    \begin{align}
        \phi &= \sum(x^{T}\mathbf{Q}x+u^{T}\mathbf{R}u) \\
        \mathbf{Q} &= \begin{bmatrix}
        1 & 0 & 0 & 0 & 0 & 0         \\
        0 & 0 & 0 & 0 & 0 & 0         \\
        0 & 0 & 0 & 0 & 0 & 0         \\
        0 & 0 & 0 & 0 & 0 & 0         \\
        0 & 0 & 0 & 0 & 0 & 0         \\
        0 & 0 & 0 & 0 & 0 & 0         
        \end{bmatrix}\\
        \mathbf{R} &= \begin{bmatrix}
            q_1        & 0        \\
            0          & q_2        
        \end{bmatrix}
    \end{align}
\end{subequations}
Having the cost function on this form makes it very easy to change or add weighting for different states and inputs.

We restricted the helicopter to obey the non-linear inequality constraint
\begin{equation}
e_k \geq \alpha e^{-\beta(\lambda_k-\lambda_t)^{2}}
\end{equation}
every timestep $k$ until the time horizon.
Since this constraint is non-linear, a non-linear solver was needed. The matlab function \texttt{fmincon} was used to compute the optimal trajectory through the state space.

\todo[inline]{MAYBE ADD EXPLAIN FMINCON SOLUTION METHOD}

\todo[inline]{ADD PLOTS OF THE OPTIMAL PATH!!}
\todo[inline]{DISCUSS RESULTS!!}

\subsection{Experimental results}
The optimal trajectory was fed into the helicopter, both with and without feedback

\todo[inline]{ADD PLOTS OF THE EXPERIMENT PATH!!}
\todo[inline]{DISCUSS RESULTS!!}


\end{document}
